{\small
\begin{longtable}{@{}>{\raggedright}p{0.4cm}>{\raggedright}p{2.5cm}>{\raggedright}p{4.2cm}>{\raggedright}p{1.5cm}p{6.2cm}@{}}
\caption{Data-quality problems found in the real samples (generated by the profiling notebook).}\label{tab:dq}\\
\toprule \# & Source & Problem & Affected & Example record / design impact\\ \midrule \endhead
1 & \texttt{noaa\_\allowbreak storm\_\allowbreak events} & Local timestamps without offset; zone stored apart in CZ\_TIMEZONE with several labels (CST-6, MST-7); two-digit year & 724/724 & \textit{EVENT\_ID 1256799: BEGIN\_DATE\_TIME='21-APR-25 10:30:00', CZ\_TIMEZONE='CST-6'}. Convert to UTC using the offset in CZ\_TIMEZONE before any comparison with other sources.\\
2 & \texttt{noaa\_\allowbreak storm\_\allowbreak events} & Damage stored as text with K/M/B suffix; empty values mixed with '0.00K' & 724/724 & \textit{EVENT\_ID 1256799: DAMAGE\_PROPERTY='0.00K' -> 0 USD}. Parse to numeric USD; keep empty as NULL (unknown), not 0.\\
3 & \texttt{noaa\_\allowbreak storm\_\allowbreak events} & No physical severity measure for floods (MAGNITUDE empty); severity must be derived & 724/724 & \textit{EVENT\_ID 1256799 (Flood): MAGNITUDE empty}. Severity = rainfall from GHCN or impact (deaths/damage), not a native field.\\
4 & \texttt{ncei\_\allowbreak ghcnd} & Trace precipitation coded as PRCP=0 with measurement flag 'T' & 232/1825 & \textit{USW00003927 2025-01-05: PRCP=0.0, attributes='T,,W'}. Keep the flag in the domain table; 0 mm with T is not the same as dry.\\
5 & \texttt{ncei\_\allowbreak ghcnd} & Daily value is a local-day total with no time zone; observation time in OBSTIME (values seen: (empty), 2400) & 1825/1825 & \textit{USW00003927 2025-01-01: OBSTIME='(empty)', attributes=',,W'}. Match to events by local calendar day (+/- 1 day), not by exact UTC instant.\\
6 & \texttt{nws\_\allowbreak observations} & precipitationLastHour is null in many observations (value: null, unitCode wmoUnit:mm) & 2463/2500 & \textit{KAUS 2026-11-06T02:55:00+00:00: precipitationLastHour=null, qc=nan}. Only the hourly (METAR) reports carry precipitation; sum precipitationLastHour from those, never from 5-min reports.\\
7 & \texttt{nws\_\allowbreak observations} & Mixed reporting intervals, not hourly (median 5 min; most common: 5 min x954, 15 min x205, 10 min x151) & 2461/2500 & \textit{KAUS 2026-11-01T13:40:00+00:00: 5 min after previous observation}. Aggregate to fixed hourly/daily buckets; watermark per station on timestamp, not on count.\\
8 & \texttt{nws\_\allowbreak alerts} & Alerts without geometry (only zone/county codes in affectedZones, text in areaDesc) & 89/226 & \textit{a938581ca63053af00abcf6f715cdb.001.1 Dense Fog Advisory: geometry=null, areaDesc='Caddo; Bossier; Webster; Claiborne; Lincoln; Union; De Soto;'}. Need NWS zone/county geometries to map alerts; join on UGC/SAME codes.\\
9 & \texttt{nws\_\allowbreak alerts} & Timestamps in local offsets, not UTC (offsets seen: -05:00, -06:00, -07:00) & 226/226 & \textit{324c74ecf2274302ec477fdf42445c.001.1: sent='2026-11-05T20:40:00-06:00'}. Normalize every timestamp to UTC on ingestion.\\
10 & \texttt{nws\_\allowbreak alerts} & Same warning re-issued as Update/Cancel with a new id (chain via references) & 111/226 & \textit{50cad2923b9d7c15f9df1b37f76079.001.1 Update Dense Fog Advisory: 1 reference(s)}. Group alerts by reference chain; the latest message replaces earlier ones.\\
11 & \texttt{nws\_\allowbreak alerts} & Feed mixes all hazards; flood-related events must be filtered by 'event' & 191/226 & \textit{7dae43678f154cccbf5d8510819d7d.001.1: event='Special Weather Statement'}. Filter on the event list in config/sources.yaml.\\
\bottomrule
\end{longtable}
}